\section{Introduction}

Cognitive science suggests that humans typically require multiple rounds of thinking to arrive at correct conclusions, especially when dealing with complex problems~\citep{nelson1990metamemory}. Inspired by this principle, recent multi-round reasoning techniques, such as Chain-of-Thought~\citep{wei2022chain}, debating~\citep{pmlr-v235-khan24a}, and self-refinement~\citep{madaan2024self,liu2024large}, applied to auto-regressive Large Language Models (LLMs), have achieved significant success across various reasoning tasks, including mathematical problem solving~\citep{cobbe2021training}, commonsense reasoning~\citep{talmor2019commonsenseqa}, scientific question answering~\citep{clark2018think}, and multi-hop question answering~\citep{yang2018hotpotqa}. This exceptional capability is widely attributed to the in-context learning abilities of auto-regressive language models~\citep{li2024dissecting,yang2024context}.

In-context learning ability has led to the conjecture that auto-regressive generative models can simulate Turing machines~\citep{lichain,merrill2024the}. \citet{schuurmans2023memory} has shown that large language models, particularly when augmented with memory, can execute complex computational processes that resemble the behavior of Turing machines. \citet{pmlr-v235-malach24a} has explored the theoretical underpinnings of auto-regressive models and their connection to universal computation, showing that even simple next-token predictors can, under the right conditions, emulate the behavior of Turing machines. \citet{lichain} illustrates the solvability of the problems that belong to $\texttt{AC}^0$, a proper subset of $\texttt{NC}^0$, via Chain-of-Thought. Yet, %The bad part is that 
the existence of some strong assumptions of these works, such as the dependence on external memory~\citep{schuurmans2023memory}, the presence of Chain-of-Thought in the training data~\citep{pmlr-v235-malach24a}, or the infinite number of layers~\citep{yuntransformers}, do not explain well the actual working conditions of realistic multi-round language models.

To make matters worse, the computability perspective does not provide a plausible explanation for how machine learning models learn. This is because learning ability and reasoning ability are not inherent properties of a computable class. An extreme example is that even games like Magic: The Gathering and Super Mario can be Turing complete~\citep{churchill2020magic}. %\dcheng{This sounds too extreme. Maybe: However, it is quite implausible that we would build a generalized AI...--} 
But we're not going to build a generalized AI out of a card game. Beyond the approximation of Turing machines, we need to understand (1) the ability of multi-round auto-regressive language models to approximate functions, (2) whether such an ability is learnable, and what its learning complexity is, and (3) the ability to generalize when inferring with imperfectly-trained language models in reality.

In this work, we systematically investigate why multi-round reasoning can improve the overall large language model problem-solving capability. Particularly, we study the approximation, learnability, and generalization abilities of auto-regressive generative models with a limited context window. Unlike previous theories, this series of studies corresponds closely to real-world scenarios, providing empirical guidance on training and inference auto-regressive generative language models.

%\chenhui{This paragraph need rewrite, the claim is not true.} For the approximation ability of the model, we show that a problem for which there exists a solution that a finite language can describe, then it can be approximated by a finite language model in multiple rounds of reasoning under the auto-regressive paradigm.
We begin with the approximation ability of auto-regressive Transformers. For the approximation ability, we show that Transformers with finite context window size are universal approximators for some steps of a Turing-computable function. Further, we prove that any Turing-computable sequence-to-sequence function can be approximated by a multi-round auto-regressive generation process. This demonstrates the feasibility of the current dominant language models for solving problems. It is worth noting that although this finding does not directly explain the problem-solving ability of language models, it is a cornerstone of the generalizability of auto-regressive models.

Next, we turn our attention to learnability, which is a critical aspect of understanding how auto-regressive language models, particularly in a multi-round reasoning context, gain their problem-solving capabilities. We first expand probably approximately correct (PAC) learning~\citep{valiant1984theory} to finite-size window next token prediction, {\color{black}and to auto-regressive sequence generation}. Beyond that, we generalize the finite-window sequence learnability to the case of exceeding the window size by means of multi-round language generation. The results show that even if the required sequence length exceeds the maximum window of the auto-regressive language model, the model remains learnable for long sequence complex problems. The sample complexity required to learn the ability to auto-regressively predict an entire long sequence will increase dramatically compared to simply making a single-word prediction. Further, we show that training with the multi-round generation paradigm has an exponential impact on the sample complexity w.r.t the number of rounds {\color{black}\(R\)}. %This means that the main role of the multi-round generation paradigm via autoregressive language modeling is not in the training phase, but in the inference phase.

Then, we focus on the generalization ability through multi-round reasoning. In particular, we show that the generalization error of the model grows with the propagation rounds. Through our analysis, we can conclude that as the number of rounds of model generation increases, eventually the answers it obtains will diverge. Nevertheless, we can still constrain the state of the intermediate process, e.g., by providing some hints in the multi-round dialogues, to control the generalization of the model and induce it to the answer we want. We point out that prompting tricks like Chain-of-Thought, self-refinement, and multiple rounds of dialogue serve to constrain the generalization error during the inference process so that the answers generated are within our expectations.

\textbf{Contributions.} Overall, we make the following contributions:
\begin{itemize}
    \item We comprehensively investigate the approximation, learning, and generalization capabilities of finite context auto-regressive language models from a theoretical perspective.
    \item We theoretically identify a dramatic increase in the sample complexity required to induce a finite window next-word prediction into sequence generation and remark that this increase can be mitigated by introducing multiple rounds of sub-sequence generation.
    \item We theoretically analyze the inter-round propagation mechanism of the generalization error during the multi-round generation process. We also point out that without intervention in the generation of multiple rounds, their cumulative error will not guarantee controllability. %We find that it is effective to intervene through methods such as CoT.
\end{itemize}

The remaining paper is structured as follows: In Section~\ref{section2}, we discuss related research on multi-round language generation, language model learning, and generalization capabilities, followed by an introduction to some relevant lemmas and definitions in Section~\ref{section3}. We theoretically analyze the approximation capabilities of multi-round Transformers in Section~\ref{section4}, the learnability of language generation with auto-regressive language models in Section~\ref{section5}, and the error propagation of generalization in Section~\ref{section6}. We conclude and provide some insights and future directions in Section~\ref{section7}.